% Conclusiones RELEVANTES del estudio realizado. Deber\'ia responder a las preguntas: ?`todas las t\'ecnicas resuelven el mismo problema o hay algunas diferencias?, ?`En qu\'e se parecen o difieren las t\'ecnicas en el contexto del problema?, ?`qu\'e limitaciones tienen?, ?`qu\'e t\'ecnicas o estrategias son las m\'as prometedoras?, ?`existe trabajo futuro por realizar?, ?`qu\'e ideas usted propone como lineamientos para continuar con investigaciones futuras?

%Se pudo observar que para el problema de TDPDP se han realizado tanto enfoques de técnicas completas como de heurísticas y metaheurísticas. En un $53\%$ de los casos se utilizaron técnicas heurísticas y metaheurísticas para resolver los problemas, partiendo de soluciones iniciales obtenidas con métodos completos. Esto se debe a la naturaleza NP-Hard del problema, la cual fue evidenciada dentro del desarrollo del informe.

%Con base en las soluciones, se evidenció la diferencia en el enfoque de cómo resolver el problema. Dado que el problema es un VRP modificado, se observó que una de las técnicas más utilizadas fue la de resolver primeramente un viaje para el camión y luego ir revisando en qué secciones se podría cambiar la entrega por la de un dron. También resultó útil diferenciar los problemas según la cantidad de drones y camiones, para así enfocarse en solucionar una parte del problema que sea menos compleja que el total.

%La mayoría de las limitaciones se pueden resumir en la complejidad del problema. Dado que es NP-Hard, se tarda mucho tiempo en obtener una solución del problema que no se vea afectada por largos tiempos de computación. Esto es debido al tamaño del espacio de búsqueda y a las restricciones de sincronizar un dron con un camión. Un factor que se observó fue cómo las restricciones van aumentando mientras el problema se vuelve más cercano a la realidad. En los primeros modelos, las restricciones eran muy ideales y no consideraban tanto la realidad del problema. Trabajos posteriores arreglan esto, agregando gasto de batería variable según el peso del dron, o permitiendo que el dron se una a mitad de camino con un camión, entre otros.

%Las tendencias muestran que el trabajo se encuentra focalizado en hacer uso de técnicas mixtas, que combinan el modelo matemático puro con las heurísticas para la resolución de los nodos utilizados para el dron. El trabajo futuro que queda por realizar sería hacer un análisis más profundo de las diferencias entre los modelos (según cantidad de drones, camiones, batería) y cómo estos modelos se relacionan con la realidad de la logística, para que así este trabajo pueda desarrollar un mayor valor en las empresas de entrega.

El trabajo realizado sobre el problema de Truck Drone Parcel Delivery Problem (TDPDP) logro mostrar la complejidad de encontrar soluciones buenas para su naturaleza NP-HARD. 
Se mostraron diversas variantes en la literatura vistas en \ref{tab:variantes-tpdpd}, donde se puede variar como enfocar el dron, utilizar distintos camiones y drones, o incluso hacer posible que el dron se encuentre a mitad de camino con el camión.


Para obtener soluciones del TPDPD, se han realizado tanto enfoques de técnicas completas como de heurísticas y metaheurísticas, en donde el $53\%$ de los casos apostaba por técnicas incompletas
Mediante la revisión del estado de Arte, se pudo observar como se han utilizados varios enfoques basados en diversas técnicas incompletas como Simmulated Annealing, Tabu Search y Algoritmos Evolutivos. 
El enfoque que se le encontro mayor prioridad fue el Truck-First, Drone-Second, en donde se trabaja el problema como un TSP y se realiza un algoritmo SPLIT que transforma al TSP en TDPDP.

Ante este panorama, se decidio hacer un análisis de una implementación evolutiva hibrida con el enfoque Truck-First, Drone-Second. 
El algoritmo evolutivo trabaja sobre la secuencia TSP que luego se transforma en TDPDP mediante una asignación SPLIT que utiliza programación dinámica. 
Además, el algoritmo luego es mejorado mediante una generación de individuos de mayor calidad mediante un GRASP que optimiza a uno de los mejores $k=5$ vecinos.
Tanto la mutuación como el cruzamiento fortalecen los puntos fuertes del TSP, sin tener que preocuparse de las restricciones del TDPDP las cuales son tratadas por el DP. 
Especificamente, en mutuación se implementa una estrategia adaptativa que permite mejorar los individuos según su calidad relativa al mejor individuo actual.

Con esto, se realiza una experimentación que permita ver el impacto de los parámetros. 
Se utilizaron tanto parametros sintonizados mediante paramILS como parametros especificos y casos bordes para poder comprender su impacto en la evaluación y tiempo de la ejecución de l algoritmo.

En instancias pequeñas de $c \in \{10,50\}$ se alcanzaron los valores optimos, pero con $c \in \{100,200\}$ se obtuvieron valores que aún podian ser mejorados. 
Con esto se pudo observar el aporte de utilizar GRASP para individuos iniciales de mejor caldad, con una mejora incluso de $40.000$ en comparación la evaluación de los algoritmos que utilizaron individuos iniciales aleatorios.

Se pudo observar también como que configuraciones extremas con poca mutuación o poca exploración dan peores valores que aquellas mas equilibradas. El punto optimo fue $p_{mut} = 0.2$ y $p_{cro} = 0.3$
También se observo como bajo un $n \times g$ identico, las soluciones entregadas por los algoritmos tienen un tiempo de ejecución parecido, pero con mejores soluciones cuando $n$ tiende a ser más pequeño que $g$.

Para un futuro se propone hacer un análisis más profundo sobre las configuraciones posibles, considerando más casos en el medio que casos bordes. También se propone hacer una comparación con implementaciones evolutivas que si consideren al dron en su mutuación y cruzamiiento, para observar como se enfrenta con el enfoque Truck-First, Drone-Second, el cual sigue tiendo la limitación de la miopía en su desarrollo
